\documentclass[11pt]{article}
\usepackage[margin=1in]{geometry}
\usepackage[T1]{fontenc}
\usepackage[utf8]{inputenc}
\usepackage{lmodern}
\usepackage{amsmath, amssymb}
\usepackage{booktabs}
\usepackage{hyperref}
\usepackage{enumitem}
\usepackage{graphicx}

\setlist[itemize]{leftmargin=1.5em}
\setlist[enumerate]{leftmargin=1.8em}

\newcommand{\placeholder}[1]{\textit{[#1]}}

\title{CS336 Assignment 1 (Basics)\\Writeup}
\author{Lenard Strahringer}
\date{\today}

\begin{document}
\maketitle
\tableofcontents
\newpage

\section{Problem (unicode1): Understanding Unicode}
\subsection*{(a)}
\texttt{chr(0)} returns the null character (U+0000), a not visible control character.

\subsection*{(b)}
Its string representation shows \texttt{'\textbackslash x00'}, while printing it produces a non-printing control character (visually empty output).

\subsection*{(c)}
When this character appears in text it is embedded as an invisible null byte that still occupies a position in the string.

\section{Problem (unicode2): Unicode Encodings}
\subsection*{(a)}
UTF-8 is usually preferred. It is much more space-efficient for ASCII-heavy text (common on the web). It is byte-oriented and backward-compatible with ASCII. It also avoids UTF-16 endianness and BOM complications while still representing all Unicode characters.

\subsection*{(b)}
An example is \texttt{b'\textbackslash xc3\textbackslash xa9'} (which is ``é'' in UTF-8). The function is incorrect because it decodes each byte independently. UTF-8 is variable-length, and many characters require multiple bytes for encoding.

\subsection*{(c)}
An example invalid two-byte sequence is \texttt{b'\textbackslash xc0\textbackslash xaf'}, which is an overlong (i.e., it uses more bytes than needed to represent a character) UTF-8 encoding and therefore does not decode to any valid character.

\section{Problem (train\_bpe\_tinystories): BPE Training on TinyStories}
\subsection*{(a)}
Using \texttt{input\_path=data/TinyStoriesV2-GPT4-train.txt}, \texttt{max\_vocab\_size=10000}, training took 137.01 seconds with a peak RAM usage of 175.2\,MB. The resulting vocabulary had a longest token of 15 bytes: \texttt{' accomplishment'}. This makes sense because BPE tends to merge frequent whitespace-prefixed word fragments into single tokens.

\subsection*{(b)}
In my training process, most of the runtime is in the BPE merge loop. \texttt{update\_word\_pairs} (32.38\%) and \texttt{merge\_tokens\_words\_element} (32.16\%) together account for about two-thirds of sampled execution time. Smaller portions come from loop overhead and preprocessing.

\section{Problem (train\_bpe\_expts\_owt): BPE Training on OpenWebText}
\subsection*{(a)}
Using \texttt{input\_path=data/owt\_valid.txt}, \texttt{max\_vocab\_size=32000}, training took 1943.20 seconds. It produced a 32,000-token vocabulary whose longest token was 64 bytes: \texttt{'----------------------------------------------------------------'}. This token is plausible because long repeated punctuation spans occur in web text and are highly mergeable under BPE.

\subsection*{(b)}
The OpenWebText tokenizer (32k vocab) produces noticeably different merge patterns from the TinyStories tokenizer (10k vocab). The OWT tokenizer learns longer, more varied tokens reflecting web-text diversity, URL fragments, and technical vocabulary. The TinyStories tokenizer is dominated by common narrative words and whitespace-prefixed word fragments typical of simple children's stories. As a result, the OWT tokenizer achieves better compression on web text. The TinyStories tokenizer is better tuned to the narrow vocabulary of the stories domain, but it also includes fewer full words as tokens (and many subparts instead).

\section{Problem (tokenizer\_experiments): Experiments with Tokenizers}
\subsection*{(a)}
On 10 sampled TinyStories documents, the TinyStories tokenizer achieved 4.150 bytes/token (7412 bytes over 1786 tokens). On 10 sampled OpenWebText documents, the OpenWebText tokenizer achieved 4.680 bytes/token (31487 bytes over 6728 tokens). The OWT tokenizer, therefore, achieves better compression on its sample.

\subsection*{(b)}
Tokenizing the OpenWebText sample with the TinyStories tokenizer produced 3.160 bytes/token (31487 bytes over 9963 tokens). The OpenWebText tokenizer on the same sample produced 4.680 bytes/token. The TinyStories tokenizer produced many more tokens (worse compression).

\subsection*{(c)}
Measured throughput was about 4.12 MB/s for the TinyStories tokenizer and 1.92 MB/s for the OpenWebText tokenizer (benchmarked on about 5.0M bytes each). At these rates, tokenizing the 825 GB Pile would take about 55.67 hours (2.32 days) with the TinyStories tokenizer. With the OpenWebText tokenizer, it would take about 119.08 hours (4.96 days).

\subsection*{(d)}
\texttt{uint16} is appropriate because both vocabularies are below 65,536 entries (10K and 32K), so every token ID fits exactly in 16 bits. Using \texttt{uint16} therefore halves storage relative to \texttt{int32} while preserving all IDs.

\section{Problem (transformer\_accounting): Transformer LM Resource Accounting}
\subsection*{(a)}
Trainable parameters:
\[
\text{total}
= (50{,}257\cdot1{,}600)
+ 48\left(4\cdot1{,}600^2 + 3\cdot1{,}600\cdot4{,}288 + 2\cdot1{,}600\right)
+ 1{,}600
+ (1{,}600\cdot50{,}257)
\]
\[
= 80{,}411{,}200 + 1{,}479{,}628{,}800 + 1{,}600 + 80{,}411{,}200
= 1{,}640{,}452{,}800
\]

If each parameter is stored as a single-precision float, the model needs $32 \times 1{,}640{,}452{,}800 = 52{,}494{,}489{,}600\text{ bits} \approx 6.56\text{ GB}$. At least 6.56~GB of memory is required just to load the model.

\subsection*{(b)}
For one \texttt{TransformerBlock}, let \(B=\) batch size, \(T=\) sequence length, \(d=1600\), \(h=25\), \(d_k=d/h=64\), \(d_{\text{ff}}=4288\). The 9 matrix multiplications are:
\begin{itemize}
    \item Q projection: \([B,T,d]\times[d,d]\to[B,T,d]\)
    \item K projection: \([B,T,d]\times[d,d]\to[B,T,d]\)
    \item V projection: \([B,T,d]\times[d,d]\to[B,T,d]\)
    \item Attention scores \(QK^\top\): \([B,h,T,d_k]\times[B,h,d_k,T]\to[B,h,T,T]\)
    \item Attention-value product: \([B,h,T,T]\times[B,h,T,d_k]\to[B,h,T,d_k]\)
    \item Output projection \(W_O\): \([B,T,d]\times[d,d]\to[B,T,d]\)
    \item FFN branch \(W_1\): \([B,T,d]\times[d,d_{\text{ff}}]\to[B,T,d_{\text{ff}}]\)
    \item FFN branch \(W_3\): \([B,T,d]\times[d,d_{\text{ff}}]\to[B,T,d_{\text{ff}}]\)
    \item FFN output \(W_2\): \([B,T,d_{\text{ff}}]\times[d_{\text{ff}},d]\to[B,T,d]\)
\end{itemize}

For \texttt{TransformerLM} outside the \texttt{TransformerBlock}s:
\begin{itemize}
    \item Final output projection: \([B,T,d]\times[d,\text{vocab\_size}]\to[B,T,\text{vocab\_size}]\)
\end{itemize}

One forward pass needs \(9\times48+1=433\) matrix multiplications.

General rule for matrix multiplication FLOPs:
\[
([m,n]\times[n,p]) \Rightarrow 2mnp \text{ FLOPs}
\]

Using \(B=1\), \(T=1024\), \(d=1600\), \(h=25\), \(d_k=64\), \(d_{\text{ff}}=4288\), \(\text{vocab\_size}=50{,}257\):
\begin{itemize}
    \item Q, K, V, \(W_O\): each \(2BTdd=5{,}242{,}880{,}000\)
    \item \(QK^\top\), attention-value: each \(2BhTTd_k=3{,}355{,}443{,}200\)
    \item \(W_1, W_3, W_2\): each \(2BTdd_{\text{ff}}=14{,}050{,}918{,}400\)
\end{itemize}

Per-block FLOPs:
\[
4(5{,}242{,}880{,}000) + 2(3{,}355{,}443{,}200) + 3(14{,}050{,}918{,}400)
= 69{,}835{,}161{,}600
\]

All 48 blocks:
\[
48\times69{,}835{,}161{,}600
= 3{,}352{,}087{,}756{,}800
\]

Final output projection FLOPs:
\[
2BTd\cdot\text{vocab\_size}
= 2\cdot1\cdot1024\cdot1600\cdot50{,}257
= 164{,}682{,}137{,}600
\]

Total FLOPs in one forward pass:
\[
3{,}352{,}087{,}756{,}800 + 164{,}682{,}137{,}600
= \boxed{3{,}516{,}769{,}894{,}400}
\]

\subsection*{(c)}
The PositionwiseFeedForward requires $14{,}050{,}918{,}400 \times 3 = 42{,}152{,}755{,}200$ FLOPs. The CausalMultiHeadSelfAttention requires $3{,}355{,}443{,}200 \times 2 + 5{,}242{,}880{,}000 = 11{,}953{,}766{,}400$ FLOPs. The FFN costs more than four times the attention.

\subsection*{(d)}
Using the same accounting as above, the FLOP proportions are:

\begin{center}
\begin{tabular}{lrrrrrrr}
\toprule
Model & QKV & Attn scores & Attn value & Attn out & FFN W1/W3 & FFN W2 & Final proj \\
\midrule
GPT-2 small  & 14.91\% & 6.63\% & 6.63\% & 4.97\% & 26.51\% & 13.25\% & 27.10\% \\
GPT-2 medium & 18.62\% & 6.21\% & 6.21\% & 6.21\% & 33.37\% & 16.68\% & 12.70\% \\
GPT-2 large  & 20.49\% & 5.46\% & 5.46\% & 6.83\% & 36.20\% & 18.10\% & 7.45\% \\
GPT-2 XL     & 21.47\% & 4.58\% & 4.58\% & 7.16\% & 38.36\% & 19.18\% & 4.68\% \\
\bottomrule
\end{tabular}
\end{center}

As model size increases, FFN and projection-heavy terms take a larger share of total FLOPs. The final vocab projection takes a much smaller share. The attention score and value terms decrease proportionally because \(T\) is fixed, and depth-dependent parts grow faster.

\subsection*{(e)}
For GPT-2 XL, increasing context length from \(1024\) to \(16{,}384\) raises one-forward-pass FLOPs from \(3{,}516{,}769{,}894{,}400\) to \(133{,}577{,}729{,}638{,}400\) ($\approx $\(37.98\times\) higher). The FLOP share shifts strongly toward the \(T^2\) attention terms. Attention scores and attention-value rise from \(9.16\%\) to \(61.73\%\). FFN (\(W_1/W_3/W_2\)) falls from \(57.51\%\) to \(24.24\%\). The final projection falls from \(4.68\%\) to \(1.97\%\).

\section{Problem (learning\_rate\_tuning): Tuning the Learning Rate (SGD toy example)}
At \(\eta=10^1\), the loss decreases steadily but slowly over 10 iterations. At \(\eta=10^2\), the loss decreases faster and reaches a lower final value. This indicates a more aggressive but still stable step. At \(\eta=10^3\), the loss diverges and increases over the course of training. This learning rate is far too large for the quadratic toy loss landscape.

\section{Problem (adamw\_accounting): Resource Accounting for Training with AdamW}
\subsection*{(a)}
Let \(B=\text{batch\_size}\), \(V=\text{vocab\_size}\), \(T=\text{context\_length}\), \(L=\text{num\_layers}\), \(d=d_{\text{model}}\), \(h=\text{num\_heads}\), \(d_{\text{ff}}=\frac{8}{3}d\). Float32 uses 4 bytes/element.

\[
P = 2Vd + L\!\left(4d^2 + 3dd_{\text{ff}} + 2d\right) + d
= 2Vd + L(12d^2 + 2d) + d .
\]

\[
M_{\text{params}} = 4P \text{ bytes}.
\]


\[
\begin{aligned}
A_{\text{block}} &= \underbrace{2BTd}_{\text{2 RMSNorms}}
+ \underbrace{(3BTd + BhT^2 + BhT^2 + BTd + BTd)}_{\text{QKV, }QK^\top,\text{ softmax, weighted sum, }W_O}
+ \underbrace{(BTd_{\text{ff}} + BTd_{\text{ff}} + BTd_{\text{ff}} + BTd_{\text{ff}} + BTd)}_{\text{W1, W3, SiLU, product, W2}} \\
&= 8BTd + 4BTd_{\text{ff}} + 2BhT^2
= \frac{56}{3}BTd + 2BhT^2 .
\end{aligned}
\]

\[
A_{\text{extra}} = BTd + 2BTV.
\]

\[
A = L A_{\text{block}} + A_{\text{extra}}
= L\!\left(\frac{56}{3}BTd + 2BhT^2\right) + BTd + 2BTV .
\]

\[
M_{\text{act}} = 4A \text{ bytes}.
\]


\[
M_{\text{grad}} = 4P \text{ bytes}.
\]


\[
M_{\text{opt}} = 8P \text{ bytes}.
\]


\[
M_{\text{peak}} = M_{\text{params}} + M_{\text{act}} + M_{\text{grad}} + M_{\text{opt}}
= 16P + 4A \text{ bytes},
\]
\[
\boxed{
M_{\text{peak}}
= 32\,\text{vocab\_size}\,d_{\text{model}}
+ 192\,\text{num\_layers}\,d_{\text{model}}^2
+ 32\,\text{num\_layers}\,d_{\text{model}}
+ 16\,d_{\text{model}}
+ \frac{224}{3}\,\text{num\_layers}\,\text{batch\_size}\,\text{context\_length}\,d_{\text{model}}
+ 
8\,\text{num\_layers}\,\text{batch\_size}\,\text{num\_heads}\,\text{context\_length}^2
+ 4\,\text{batch\_size}\,\text{context\_length}\,d_{\text{model}}
+ 8\,\text{batch\_size}\,\text{context\_length}\,\text{vocab\_size}
\;\text{bytes}.}
\]

\subsection*{(b)}
Using \(V=50{,}257\), \(T=1024\), \(L=48\), \(d=1600\), \(h=25\), \(d_{\text{ff}}=4288\):
\[
M_{\text{peak}}(B)=aB+b
=16{,}373{,}391{,}360\cdot B + 26{,}247{,}244{,}800
\quad\text{bytes}.
\]
Under an 80GB budget (\(80{,}000{,}000{,}000\) bytes), the largest integer batch size is
\[
B_{\max}=\left\lfloor\frac{80{,}000{,}000{,}000-26{,}247{,}244{,}800}{16{,}373{,}391{,}360}\right\rfloor=3.
\]
Check: \(M_{\text{peak}}(3)=75{,}367{,}418{,}880<80{,}000{,}000{,}000\), while \(M_{\text{peak}}(4)=91{,}740{,}810{,}240>80{,}000{,}000{,}000\).

\subsection*{(c)}
\begin{align*}
    \text{1. Weight decay:} \quad & p \gets p - (\eta \lambda) p && \text{(2 FLOPs)} \\
    \text{2. First moment ($m$):} \quad & m \gets \beta_1 m + (1 - \beta_1) g && \text{(3 FLOPs)} \\
    \text{3. Second moment ($v$):} \quad & v \gets \beta_2 v + (1 - \beta_2) g^2 && \text{(4 FLOPs)} \\
    \text{4. Parameter update:} \quad & p \gets p - \alpha_t \frac{m}{\sqrt{v} + \epsilon} && \text{(5 FLOPs)} \\[1ex]
    \textbf{Total:} \quad & \textbf{14 FLOPs per parameter} \implies \mathbf{14P \text{ FLOPs}} &&
\end{align*}

\subsection*{(d)}
\[
202.086 \text{ days} \times 24 \text{ hours/day} \approx \mathbf{4,850 \text{ hours}}.
\]
One forward pass for one sequence costs \(3{,}516{,}769{,}894{,}400\) FLOPs (from the transformer\_accounting problem). Multiplying by \texttt{batch\_size} gives forward FLOPs per step. Using \(\approx 3\times\) forward FLOPs for forward+backward gives training FLOPs per step. At 50\% MFU on one H100, effective throughput is \(495\times 10^{12}\times 0.5\) FLOPs/s. Training time is total FLOPs divided by this rate.

\section{Problem (training\_together): Put it Together}
The script for Question 5.3 (training loop) is in \texttt{cs336\_basics/train\_lm.py}.

\section{Problem (decoding): Decoding}
The script for Question 6.1 (decoding/generation) is in \texttt{cs336\_basics/train\_lm.py}.

\section{Problem (experiment\_log): Experiment Logging}
For Question 7.1, experiment logging is implemented in \texttt{cs336\_basics/run\_logged\_experiment.py} and \texttt{cs336\_basics/plot\_experiment\_log.py}, with core train-loop logging in \texttt{cs336\_basics/train\_lm.py}.

\section{Problem (learning\_rate): Tune the Learning Rate on TinyStories}
\subsection*{(a)}
I ran a log-scale learning-rate sweep with fixed \texttt{batch\_size}=128 and \texttt{max\_iters}=10{,}000. The sweep covered
\[
\eta \in \{10^{-4}, 3\times 10^{-4}, 10^{-3}, 3\times 10^{-3}, 10^{-2}, 3\times 10^{-2}, 10^{-1}, 1\}.
\]
All runs therefore used the same token budget:
\[
\text{batch\_size}\times \text{max\_iters}\times \text{context\_length}=327{,}680{,}000,
\]
and results were compared using validation loss from the logged checkpoints.

Best observed \texttt{val\_loss} for each learning rate:
\begin{center}
\begin{tabular}{lrr}
\toprule
Learning rate \(\eta\) & Best \texttt{val\_loss} & Final \texttt{val\_loss} \\
\midrule
\(10^{-4}\) & 1.7587 & 1.7641 \\
\(3\times 10^{-4}\) & 1.5148 & 1.5321 \\
\(10^{-3}\) & 1.3784 & 1.3784 \\
\(3\times 10^{-3}\) & 1.3229 & 1.3278 \\
\(10^{-2}\) & \textbf{1.3107} & \textbf{1.3107} \\
\(3\times 10^{-2}\) & 1.3460 & 1.3493 \\
\(10^{-1}\) & 1.9483 & 1.9807 \\
\(1\) & 2.8287 & 2.8287 \\
\bottomrule
\end{tabular}
\end{center}

Best run reached \(\texttt{val\_loss}=1.3107\) at \(\eta=10^{-2}\).

\subsection*{(b)}
None of these runs showed hard numerical divergence (no NaNs/Infs, all runs reached step 10{,}000). However, optimization quality varied strongly with \(\eta\). The \(\eta=1\) curve shows large early instability spikes in validation loss, then recovers to a much worse plateau. The best region is \(\eta \in [3\times 10^{-3}, 10^{-2}]\). Too-small rates (\(10^{-4}, 3\times 10^{-4}\)) underfit within the token budget. Too-large rates (\(10^{-1}, 1\)) train to clearly worse final loss.

\begin{figure}[htbp]
    \centering
    \includegraphics[width=0.99\textwidth]{\detokenize{figures/learning_rate/lr_sweep_20260424.png}}
    \caption{Validation-loss curves from the learning-rate sweep submitted in this run (\(\eta \in \{10^{-4},3\times10^{-4},10^{-3},3\times10^{-3},10^{-2},3\times10^{-2},10^{-1},1\}\), \texttt{batch\_size}=128). The \(\eta=1\) trace is highly unstable early, while mid-range rates achieve the best final losses.}
    \label{fig:lr-divergence}
\end{figure}

The best rate (\(10^{-2}\)) is approximately \(1/10\) of the rate where instability starts (\(10^{-1}\)).


\section{Problem (batch\_size\_experiment): Batch Size Variations}
Using the same W\&B export, I compared the best retuned learning rate for each batch size:
\begin{figure}[htbp]
    \centering
    \includegraphics[width=0.49\textwidth]{\detokenize{figures/batch_size/learning_curves_incl_b1.png}}
    \hfill
    \includegraphics[width=0.49\textwidth]{\detokenize{figures/batch_size/learning_curves_excl_b1.png}}
    \caption{Batch-size learning curves at each batch's best-tuned learning rate. Left: including \texttt{batch\_size}=1. Right: excluding \texttt{batch\_size}=1 for readability of practical batch sizes.}
    \label{fig:batch-size-best-lr-curves}
\end{figure}
\begin{center}
\begin{tabular}{lrr}
\toprule
Batch size & Best learning rate & Best \texttt{eval/train\_loss} \\
\midrule
1 & \(10^{-4}\) & 1.1619 \\
64 & \(10^{-3}\) & 1.3606 \\
128 & \(3\times 10^{-3}\) & 1.3427 \\
256 & \(10^{-2}\) & \textbf{1.3116} \\
1024 & \(10^{-2}\) & 1.3627 \\
\bottomrule
\end{tabular}
\end{center}

\texttt{batch\_size}=1 at \(\eta=10^{-4}\) has the lowest training loss (\(1.1619\)), but its validation loss is \(1.6665\), indicating overfitting. Among practical batch sizes, increasing from 64 to 256 improved the best achievable loss. Pushing to 1024 hurt performance again. The best practical setting was \texttt{batch\_size}=256 with \(\eta=10^{-2}\) (\texttt{eval/train\_loss} \(=1.3116\)). This suggests a non-monotonic batch-size tradeoff at a fixed token budget.

\section{Problem (generate): Generate Text}
The following is LM-generated text from a TinyStories run. Settings: \texttt{vocab\_size}=32000 (32k BPE tokenizer), learning rate \(3\times 10^{-3}\), checkpoint \texttt{tiny\_32k\_bs256\_lr3e-3\_10k\_wandb.pt} (step 10{,}000). Decoding used \texttt{temperature}=0.45, \texttt{top\_p}=0.85, and \texttt{max\_new\_tokens}=256.
\begin{quote}
Once upon a time, there was a little girl named Lilly. Lilly was three years old and loved to play with her toys. One day, she was playing in her room when she heard a noise coming from the closet. She was scared and ran to the closet to see what was going on.

Inside the closet, she found a small, scared kitten. It was meowing and seemed very sad. Lilly wanted to help the kitten, so she decided to take it to her mom.

When Lilly arrived, her mom was very surprised to see the kitten. She said, ``Oh, poor little kitten! What happened to you?''

Lilly replied, ``I was playing with my toys and I saw the closet and I wanted to help it.''

Her mom smiled and said, ``That's very kind of you, Lilly. You can help the kitten, but you must be very careful.''

Lilly nodded and said, ``I will, Mommy.'' She carefully picked up the kitten and put it back in the closet.

The next day, Lilly and her mom went to the closet and found the kitten. Lilly was so happy to help the kitten and she knew that she had done a good thing.

\texttt{<|endoftext|>}
\end{quote}

First, I had to increase the vocab size to achieve fluency. With 10k vocab, the model produced non-existent words. With 32k vocab, the model produces grammatically correct stories. Second, low temperature and top-$p$ values were needed to obtain fluent and coherent output.

\section{Problem (layer\_norm\_ablation): Remove RMSNorm and Train}
Runs (\(327.68\)M tokens):
\begin{itemize}
    \item Baseline, \(\eta=0.0028\): final \(1.3360\), best \(1.3261\).
    \item No-RMSNorm, \(\eta=0.0028\): final \texttt{NaN}, best finite \(2.2182\) (step 600).
    \item No-RMSNorm, \(\eta=0.0007\): final \(1.4111\), best \(1.4077\).
\end{itemize}
\begin{figure}[htbp]
    \centering
    \includegraphics[width=0.9\textwidth]{\detokenize{figures/q7/layernorm.png}}
    \caption{Q7 layer-norm ablation learning curves (\texttt{eval/val\_loss}): baseline RMSNorm vs.\ no-RMSNorm runs.}
    \label{fig:q7-layernorm}
\end{figure}
Removing RMSNorm significantly hurts stability and quality (Figure~\ref{fig:q7-layernorm}). The no-RMSNorm run at \(\eta=0.0028\) fails with NaN. Even the stable retuned run (\(\eta=0.0007\)) remains \(+0.0751\) worse than baseline in final validation loss.

\section{Problem (pre\_norm\_ablation): Implement Post-Norm and Train}
Runs:
\begin{itemize}
    \item Pre-norm baseline: final \(1.3320\), best \(1.3236\).
    \item Post-norm ablation: final \(1.3476\), best \(1.3456\).
\end{itemize}
\begin{figure}[htbp]
    \centering
    \includegraphics[width=0.9\textwidth]{\detokenize{figures/q7/prenorm.png}}
    \caption{Q7 pre-norm vs.\ post-norm learning curves (\texttt{eval/val\_loss}).}
    \label{fig:q7-prenorm}
\end{figure}
Interpretation: Figure~\ref{fig:q7-prenorm} shows both runs improve similarly, but pre-norm is consistently slightly lower. Final post-norm loss is \(+0.0157\) worse than pre-norm.

\section{Problem (no\_pos\_emb): Implement NoPE}
Runs:
\begin{itemize}
    \item RoPE baseline: final \(1.3243\), best \(1.3243\).
    \item NoPE ablation: final \(1.3836\), best \(1.3836\).
\end{itemize}
\begin{figure}[htbp]
    \centering
    \includegraphics[width=0.9\textwidth]{\detokenize{figures/q7/nopos.png}}
    \caption{Q7 positional-encoding ablation learning curves (\texttt{eval/val\_loss}): RoPE baseline vs.\ NoPE.}
    \label{fig:q7-nope}
\end{figure}
Interpretation: in Figure~\ref{fig:q7-nope}, both models learn quickly at first, but the RoPE baseline stays lower throughout training and converges better. NoPE is \(+0.0593\) worse than RoPE in final validation loss.

\section{Problem (swiglu\_ablation): SwiGLU vs.\ SiLU}
Runs:
\begin{itemize}
    \item SwiGLU baseline: final \(1.3308\), best \(1.3251\).
    \item SiLU ablation: final \(1.3383\), best \(1.3383\).
\end{itemize}
\begin{figure}[htbp]
    \centering
    \includegraphics[width=0.9\textwidth]{\detokenize{figures/q7/swiglu.png}}
    \caption{Q7 feed-forward activation comparison (\texttt{eval/val\_loss}): SwiGLU baseline vs.\ SiLU.}
    \label{fig:q7-swiglu}
\end{figure}
Interpretation: Figure~\ref{fig:q7-swiglu} shows nearly overlapping curves, with SwiGLU maintaining a small but consistent advantage late in training. SiLU finishes \(+0.0075\) worse than SwiGLU in final validation loss.

\section{Problem (main\_experiment): Experiment on OpenWebText}
The OpenWebText run (\texttt{owt\_bs256\_lr0.01\_10k\_vocab32k}) is shown below, compared against the TinyStories reference curve:

\begin{figure}[htbp]
    \centering
    \includegraphics[width=0.95\textwidth]{\detokenize{figures/main_experiment/owt_vs_tiny_curves.png}}
    \caption{Validation-loss curves for TinyStories vs.\ OpenWebText. The OpenWebText run starts at \(\approx 10.39\) and reaches \(\texttt{val\_loss}\approx 3.74\) by step 10{,}000, while TinyStories reaches \(\approx 1.31\) under the same token budget style.}
    \label{fig:owt-vs-tiny-val-loss}
\end{figure}

Compared with TinyStories, OpenWebText converges much more slowly and to a substantially higher loss (TinyStories best \(\approx 1.31\) vs.\ OpenWebText final \(\approx 3.74\)). This is expected. OpenWebText is more diverse, noisier, and less repetitive than TinyStories. The conditional next-token distribution has higher entropy and is harder for a small 4-layer model to fit. The higher OWT loss does not necessarily indicate a broken pipeline. It reflects a more difficult dataset and model-capacity regime.

Using the latest OpenWebText checkpoint (\texttt{checkpoints/owt\_bs256\_lr0.01\_10k\_vocab32k.pt}), generation with \texttt{max\_new\_tokens=256}, \texttt{temperature=0.9}, and \texttt{top\_p=0.95} produced a 263-token sample (no \texttt{<|endoftext|>} encountered before the length cap). The output is:

\begin{verbatim}
In a surprising turn of events, the billionaire played a major role in selling the shirts to Korean businessman Kim Il-sung. "Nobody was working in Chinese," he told reporters after receiving an Internet cover story that circulated online about the events. "He's not living in China."

Jennifer Laney, a partner with the Hymen business group, said she and her husband do not have clients. In fact, she insisted, she had no credit for having received their American Embassy notice. "It's a bad business," Laney said.

The point has changed with the notorious bitcoin vendor. In July, a Chinese hacker and a banker filed a $1 million lawsuit against 12 investors. On Thursday, an attorney at Goldman Sachs declined to comment on the case.

South Park's income in China is estimated to be about $1 billion, according to a 2011 Chinese Election Commission survey. The United States Congressmen's Office (U.S.A.) found that this year, a 20 percent increase in salaries for the registration and qualification requirements of a living student in the U.S. was the median for the state's student salary in 2013. In that same year, the country began to recover from foreign currency imports, including imports in exchange for
\end{verbatim}

The sample is locally fluent (grammatical sentences, plausible news-style phrasing), but globally inconsistent across paragraphs (topic drift, entity and claim incoherence). Two main factors affect quality. First, \textbf{decoding settings}: higher temperature and top-p increase diversity but can reduce consistency. Second, \textbf{training scale}: more high-quality tokens and a better-converged checkpoint improve long-range coherence. Additional contributors include \textbf{model capacity}, \textbf{context length}, and \textbf{tokenizer-domain match}.

\section{Problem (leaderboard): Leaderboard Submission}
Final validation loss (OpenWebText leaderboard run): \(\texttt{val\_loss}=3.743525385856629\).

\begin{figure}[htbp]
    \centering
    \includegraphics[width=0.95\textwidth]{\detokenize{figures/leaderboard/owt_val_loss_wallclock_under45m_20260424.png}}
    \caption{Leaderboard submission learning curve (\texttt{eval/val\_loss}) with wall-clock time on the x-axis (minutes). The run reaches final \(\texttt{val\_loss}=3.743525385856629\) within the 45-minute budget.}
    \label{fig:leaderboard-wallclock}
\end{figure}

I swept \texttt{batch\_size} \(\in \{64, 128, 256\}\) against \(\eta \in \{10^{-3}, 3\times10^{-3}, 10^{-2}, 3\times10^{-2}\}\) at \texttt{vocab\_size}=32000 and \texttt{context\_length}=256. \texttt{max\_iters} \(\in\{38200, 19100, 9550\}\) was paired to each batch size so \(B\cdot\texttt{max\_iters}\) (token exposure) is constant across the grid. Each job ran under a 45-minute wall-clock cap. Sweep driver: \texttt{cs336\_basics/7\_5\_owt\_batch\_lr\_sweep.py}, Modal \texttt{run\_logged\_experiment}, W\&B \texttt{experiment}=\texttt{owt\_batch\_lr\_sweep}.

Across all twelve runs, \(\eta = 10^{-2}\) achieved the lowest validation loss for every batch size. The best configuration was \texttt{batch\_size}=128, \(\eta = 10^{-2}\), \texttt{max\_iters}=19100. This used the same assignment baseline transformer: \texttt{d\_model}=512, \texttt{num\_layers}=4, \texttt{num\_heads}=16, \texttt{d\_ff}=1344, \texttt{rope\_theta}=10{,}000. AdamW settings: \(\beta_1{=}0.9\), \(\beta_2{=}0.95\), \texttt{weight\_decay}=0.1, \texttt{max\_grad\_norm}=1, \texttt{dtype}=\texttt{float32}. Additional settings: \texttt{warmup\_iters}=1910, cosine decay to \texttt{min\_learning\_rate}=\(3\times 10^{-5}\), \\ \texttt{eval\_interval}=\texttt{max\_iters}, \texttt{eval\_batches}=3, \texttt{checkpoint\_interval}=\(19100\times 10^{3}\).

That run finished in about 41 minutes. Figure~\ref{fig:owt-sweep-best-val-wall} shows \texttt{eval/val\_loss} vs.\ wall-clock time. Loss falls from \(\approx 10.4\) to \(\approx 3.7\) with an almost linear downward trend. Figure~\ref{fig:owt-sweep-best-train-wall} shows \texttt{train/loss}. It starts with a fast drop (roughly \(6.8 \rightarrow 4.7\) in the first few minutes), then a slower noisy descent toward \(\approx 3.8\). The loss was still modestly decreasing at the end, so additional steps would likely help.

\begin{figure}[htbp]
    \centering
    \includegraphics[width=0.95\textwidth]{\detokenize{figures/leaderboard/owt_sweep_best_eval_val_loss_wallclock.png}}
    \caption{Best sweep run (\texttt{bs}=128, \(\eta{=}10^{-2}\), \texttt{max\_iters}=19100): \texttt{eval/val\_loss} vs.\ wall time (minutes).}
    \label{fig:owt-sweep-best-val-wall}
\end{figure}

\begin{figure}[htbp]
    \centering
    \includegraphics[width=0.95\textwidth]{\detokenize{figures/leaderboard/owt_sweep_best_train_loss_wallclock.png}}
    \caption{Same run: \texttt{train/loss} vs.\ wall time (minutes).}
    \label{fig:owt-sweep-best-train-wall}
\end{figure}

For the separate leaderboard submission (Figure~\ref{fig:leaderboard-wallclock}), I lengthened the warmup schedule, which improved optimization stability at the higher learning rates.

\end{document}
